\documentclass[12pt,a4paper]{amsart}

\usepackage{amsmath,amssymb,amsthm}
\usepackage{mathtools}
\usepackage{hyperref}
\usepackage{geometry}
\geometry{margin=2.5cm}

% -------------------------------------------------------
% Theorem environments
% -------------------------------------------------------
\newtheorem{theorem}{Theorem}[section]
\newtheorem{lemma}[theorem]{Lemma}
\newtheorem{corollary}[theorem]{Corollary}
\newtheorem{proposition}[theorem]{Proposition}
\theoremstyle{definition}
\newtheorem{definition}[theorem]{Definition}
\newtheorem{remark}[theorem]{Remark}

% -------------------------------------------------------
% Custom commands
% -------------------------------------------------------
\newcommand{\N}{\mathbb{N}}
\newcommand{\Z}{\mathbb{Z}}
\newcommand{\Q}{\mathbb{Q}}
\newcommand{\R}{\mathbb{R}}
\newcommand{\C}{\mathbb{C}}
\newcommand{\Fp}{\mathbb{F}_p}
\DeclareMathOperator{\Re}{Re}

% -------------------------------------------------------
\begin{document}
% -------------------------------------------------------

\title{The \texorpdfstring{$L$}{L}-Function of an Elliptic Curve:\\
       Euler Product, Frobenius Trace, Analytic Continuation,\\
       and the Functional Equation}

\author{Michael Fuhrmann}
\address{Specialist Mathematics Project, Build 84}
\email{specialist-maths@localhost}
\date{\today}

\begin{abstract}
The $L$-function $L(E,s)$ of an elliptic curve $E/\Q$ encodes global
arithmetic information through its local factors.
For each prime $p$ of good reduction, the local factor is determined by
the \emph{Frobenius trace} $a_p = p + 1 - \#E(\Fp)$;
Hasse proved $|a_p| \le 2\sqrt{p}$.
The $L$-function is initially defined for $\Re(s) > 3/2$ by an Euler product
$L(E,s) = \prod_p L_p(E,s)^{-1}$ and a Dirichlet series.
The deep \emph{modularity theorem} (Wiles 1995, Breuil--Conrad--Diamond--Taylor 2001)
proves $E$ is modular, establishing analytic continuation to $\C$
and the functional equation relating $L(E,s)$ and $L(E,2-s)$.
\end{abstract}

\maketitle
\tableofcontents

% ================================================================
\section{Introduction}
% ================================================================

The $L$-function of an elliptic curve was studied by Hasse (1930s) and
Weil (1950s), motivated by the analogy with Riemann's $\zeta$ function.
The key local ingredient is the number of points $\#E(\Fp)$ for each prime $p$.
The global $L$-function packages these local data into a single analytic function,
and the Birch--Swinnerton-Dyer conjecture (Paper~27) predicts that its
behaviour at $s = 1$ determines the rank of $E(\Q)$.

The analytic continuation and functional equation were conditional on
modularity (proved by Wiles for semistable curves in 1995 \cite{Wiles1995},
completed in full generality by Breuil--Conrad--Diamond--Taylor in 2001
\cite{BCDT2001}).

% ================================================================
\section{Counting Points Modulo Primes}
% ================================================================

\begin{definition}[Frobenius trace]
\label{def:ap}
For an elliptic curve $E/\Q$ and a prime $p$ of good reduction, the
\emph{Frobenius trace} (or \emph{trace of Frobenius}) is
\[
  a_p \;=\; p + 1 - \#E(\Fp).
\]
Here $\#E(\Fp)$ counts all affine points $(x,y) \in \Fp^2$ satisfying
$y^2 = x^3 + ax + b$ (with $a, b$ reduced modulo $p$) plus the point
at infinity $\mathcal{O}$.
\end{definition}

\begin{theorem}[Hasse's theorem, 1933]
\label{thm:hasse}
For any prime $p$ of good reduction,
\[
  |a_p| \;\le\; 2\sqrt{p}.
\]
Equivalently, $\#E(\Fp) = p + 1 - a_p$ with $|a_p| \le 2\sqrt{p}$,
so $\#E(\Fp) \in [p + 1 - 2\sqrt{p},\; p + 1 + 2\sqrt{p}]$.
\end{theorem}

\begin{proof}[Proof sketch]
Write $a_p = \alpha + \bar\alpha$ where $\alpha, \bar\alpha$ are
complex numbers with $\alpha \bar\alpha = p$ (the eigenvalues of
Frobenius acting on the $\ell$-adic Tate module $T_\ell(E)$).
Then $|a_p| = |\alpha + \bar\alpha| \le 2|\alpha| = 2\sqrt{p}$.
The fact that $|\alpha|^2 = p$ is equivalent to the Riemann Hypothesis
for the zeta function of $E/\Fp$, which Hasse proved directly for
elliptic curves.
A modern proof uses the Weil cohomology theory; see \cite{Silverman1986},
Appendix~C.
\end{proof}

\begin{example}[Computing $a_p$]
\label{ex:ap}
For $E: y^2 = x^3 - x$ and $p = 5$:
The affine points modulo $5$ are those $(x,y) \in \{0,1,2,3,4\}^2$ with
$y^2 \equiv x^3 - x \pmod 5$.
\begin{center}
\begin{tabular}{c|c|c}
  $x$ & $x^3 - x \bmod 5$ & $y$ satisfying $y^2 \equiv x^3-x$ \\
  \hline
  $0$ & $0$ & $0$ \\
  $1$ & $0$ & $0$ \\
  $2$ & $6 \equiv 1$ & $1, 4$ \\
  $3$ & $24 \equiv 4$ & $2, 3$ \\
  $4$ & $60 \equiv 0$ & $0$
\end{tabular}
\end{center}
Affine points: $(0,0), (1,0), (2,1), (2,4), (3,2), (3,3), (4,0)$ --- that is $7$ points,
plus $\mathcal{O}$, giving $\#E(\mathbb{F}_5) = 8$.
Thus $a_5 = 5 + 1 - 8 = -2$.  Indeed $|-2| \le 2\sqrt{5} \approx 4.47$. \checkmark
\end{example}

% ================================================================
\section{The Euler Product}
% ================================================================

\begin{definition}[Local $L$-factors]
\label{def:local_factors}
For a prime $p$, the \emph{local $L$-factor} of $E$ at $p$ is:
\begin{itemize}
  \item \emph{Good reduction} ($p \nmid N_E$):
    \[
      L_p(E,s)^{-1} = 1 - a_p\,p^{-s} + p^{1-2s}.
    \]
  \item \emph{Multiplicative reduction} ($p \| N_E$):
    \[
      L_p(E,s)^{-1} = 1 - \epsilon_p\,p^{-s}, \quad
      \epsilon_p \in \{+1,-1\}.
    \]
  \item \emph{Additive reduction} ($p^2 \mid N_E$):
    \[
      L_p(E,s)^{-1} = 1.
    \]
\end{itemize}
\end{definition}

\begin{definition}[$L$-function of $E$]
\label{def:L_function}
The \emph{$L$-function} of $E$ is the Euler product
\[
  L(E,s) \;=\; \prod_p L_p(E,s)
  \;=\; \prod_{p \nmid N_E} \frac{1}{1 - a_p p^{-s} + p^{1-2s}}
        \cdot \prod_{p \mid N_E} \frac{1}{1 - \epsilon_p p^{-s}}.
\]
\end{definition}

\begin{proposition}[Dirichlet series representation]
\label{prop:dirichlet}
For $\Re(s) > 3/2$, the Euler product converges absolutely and equals
\[
  L(E,s) = \sum_{n=1}^{\infty} \frac{a_n}{n^s},
\]
where the coefficients $a_n$ are determined by the Euler product:
$a_1 = 1$, $a_p = $ Frobenius trace for $p \nmid N_E$,
$a_{p^k}$ determined by the recursion
$a_{p^k} = a_p\,a_{p^{k-1}} - p\,a_{p^{k-2}}$ for $p \nmid N_E$.
\end{proposition}

\begin{proof}
The Euler product converges absolutely for $\Re(s) > 3/2$ by Hasse's
bound $|a_p| \le 2\sqrt{p}$ and the estimate $|a_{p^k}| \le (k+1)p^{k/2}$.
The Dirichlet series follows from expanding each factor as a geometric series
and using multiplicativity of $n \mapsto a_n$.
\end{proof}

% ================================================================
\section{The Completed \texorpdfstring{$L$}{L}-Function and Functional Equation}
% ================================================================

\begin{definition}[Completed $L$-function]
\label{def:completed_L}
Define the \emph{conductor} $N_E$ and the \emph{completed $L$-function}
\[
  \Lambda(E,s) \;=\; \left(\frac{\sqrt{N_E}}{2\pi}\right)^s
  \Gamma(s)\,L(E,s).
\]
\end{definition}

\begin{theorem}[Functional equation via modularity]
\label{thm:functional_eq}
The completed $L$-function satisfies
\begin{equation}
  \label{eq:functional}
  \Lambda(E,s) \;=\; \epsilon(E)\,\Lambda(E,2-s),
\end{equation}
where $\epsilon(E) \in \{+1,-1\}$ is the \emph{root number} of $E$.
\end{theorem}

\begin{proof}[Proof via modularity]
By the modularity theorem
(Wiles \cite{Wiles1995} for semistable $E$,
Breuil--Conrad--Diamond--Taylor \cite{BCDT2001} in general),
$E$ corresponds to a newform $f = \sum a_n q^n$ of weight~$2$ and level~$N_E$:
for all $n$, the $n$-th Fourier coefficient of $f$ equals $a_n$.

The $L$-function of a weight-$2$ newform satisfies the functional equation
by the classical theory of modular forms (Hecke, 1936):
\[
  \Lambda(f, s) = \epsilon_f\,\Lambda(f, 2-s),
\]
where $\epsilon_f = \pm 1$ is determined by the Atkin--Lehner eigenvalue.
Since $L(E,s) = L(f,s)$, the theorem follows.
\end{proof}

\begin{remark}[Root number and parity conjecture]
The root number $\epsilon(E) = \pm 1$ has a conjectured connection to the
parity of the rank: if $\epsilon(E) = -1$, then $L(E,1) = 0$
(from the functional equation), and the BSD conjecture predicts that
the rank is odd.  The \emph{parity conjecture} asserts:
$(-1)^r = \epsilon(E)$.
\end{remark}

% ================================================================
\section{Modularity and the Connection to Modular Forms}
% ================================================================

\begin{theorem}[Modularity Theorem]
\label{thm:modularity}
Every elliptic curve $E/\Q$ is \emph{modular}: there exists a newform
$f \in S_2(\Gamma_0(N_E))$ (a cuspidal weight-$2$ eigenform for $\Gamma_0(N_E)$)
such that
\[
  L(E,s) = L(f,s) = \sum_{n=1}^{\infty} \frac{a_n(f)}{n^s},
\]
where $a_n(f)$ are the Fourier coefficients of $f$.
Equivalently, there is a non-constant morphism
$\varphi: X_0(N_E) \to E$ of algebraic curves over $\Q$.
\end{theorem}

\begin{proof}[Historical notes]
This was conjectured by Taniyama (1955) and Shimura (1957)
(the Shimura--Taniyama conjecture, later Shimura--Taniyama--Weil).
\begin{itemize}
  \item Wiles \cite{Wiles1995} proved it for semistable curves (squarefree conductor).
  \item The proof uses Galois representations attached to $E$ and to modular forms,
    and the key technical lemma: a surjection $R \to T$ (deformation ring
    to Hecke ring) is an isomorphism.
  \item Breuil, Conrad, Diamond, Taylor \cite{BCDT2001} extended the
    proof to all elliptic curves over $\Q$.
\end{itemize}
Fermat's Last Theorem follows as a corollary (via Ribet's theorem: a
hypothetical solution would give a Frey curve that is not modular).
\end{proof}

% ================================================================
\section{The \texorpdfstring{$L$}{L}-Function Near \texorpdfstring{$s=1$}{s=1}}
% ================================================================

\begin{definition}[Order of vanishing and leading coefficient]
\label{def:order_vanishing}
The \emph{order of vanishing} of $L(E,s)$ at $s=1$ is
\[
  r_{\mathrm{an}} \;=\; \ord_{s=1} L(E,s).
\]
The \emph{leading coefficient} is
\[
  L^*(E,1) \;=\; \lim_{s \to 1} (s-1)^{-r_{\mathrm{an}}} L(E,s).
\]
The BSD conjecture predicts $r_{\mathrm{an}} = r$ (the algebraic rank).
\end{definition}

\begin{remark}[Parity from the functional equation]
If $\epsilon(E) = -1$, the functional equation~\eqref{eq:functional} gives
$L(E,1) = -L(E,1)$, so $L(E,1) = 0$.
Thus $r_{\mathrm{an}} \ge 1$ whenever $\epsilon(E) = -1$.
The BSD conjecture predicts the rank is then odd.
\end{remark}

% ================================================================
\section{Numerical Evidence and Computations}
% ================================================================

\begin{example}[The curve $E: y^2 = x^3 - x$]
\label{ex:rank0}
This is the congruent number curve $E_1$ (Paper~28).
$N_E = 32$, $\epsilon(E) = +1$, $r = 0$.
$L(E,1) \approx 0.6551$ (numerically).
The BSD conjecture predicts $r = 0$ (correct, since $E(\Q) = \{O, (0,0), (\pm 1, 0)\}$,
all torsion).
\end{example}

\begin{example}[The curve $y^2 = x^3 - 25x$]
\label{ex:rank1}
The congruent number curve $E_5$ (for $n=5$).
$r = 1$, $L(E,1) = 0$, $L'(E,1) \ne 0$.
A generator of the free part is $(-4, 6)$:
$6^2 = 36$, $(-4)^3 - 25(-4) = -64 + 100 = 36$. \checkmark
\end{example}

% ================================================================
\section{Conclusion}
% ================================================================

The $L$-function of an elliptic curve is the bridge between local
arithmetic data (the Frobenius traces $a_p$) and global arithmetic
(the rank of $E(\Q)$).
Hasse's theorem bounds $|a_p| \le 2\sqrt{p}$, ensuring convergence for $\Re(s) > 3/2$.
The modularity theorem (Wiles et al.) provides analytic continuation
and the functional equation~\eqref{eq:functional} with root number $\epsilon(E) = \pm 1$.

The central question --- what is the order of vanishing $r_{\mathrm{an}}$
at $s = 1$? --- is the subject of the Birch--Swinnerton-Dyer conjecture (Paper~27).

% ================================================================
\begin{thebibliography}{10}

\bibitem{Silverman1986}
J.~H.~Silverman,
\emph{The Arithmetic of Elliptic Curves},
Graduate Texts in Mathematics~106,
Springer, New York, 1986.

\bibitem{Wiles1995}
A.~Wiles,
Modular elliptic curves and Fermat's last theorem,
\emph{Ann.\ of Math.}\ (2) \textbf{142} (1995), 443--551.

\bibitem{BCDT2001}
C.~Breuil, B.~Conrad, F.~Diamond, and R.~Taylor,
On the modularity of elliptic curves over $\Q$:
wild $3$-adic exercises,
\emph{J.\ Amer.\ Math.\ Soc.}\ \textbf{14} (2001), 843--939.

\bibitem{TaylorWiles1995}
R.~Taylor and A.~Wiles,
Ring-theoretic properties of certain Hecke algebras,
\emph{Ann.\ of Math.}\ (2) \textbf{142} (1995), 553--572.

\bibitem{Cornell1997}
G.~Cornell, J.~H.~Silverman, and G.~Stevens (eds.),
\emph{Modular Forms and Fermat's Last Theorem},
Springer, New York, 1997.

\bibitem{Knapp1992}
A.~W.~Knapp,
\emph{Elliptic Curves},
Princeton University Press, 1992.

\end{thebibliography}

\end{document}
